\section{Entropy and the covariance trace}
\label{sec:entropy}

We now turn the geometric rank estimate into the lower bound of
Theorem~\ref{ent:main}. The entropy maximizer has a smooth exponential
parameterization on each minimal face. At a hypothetical positive maximum
of the entropy deficit, the second derivative along those parameters
bounds the face dimension by a covariance trace. The rank estimate makes
that trace too small. Boundary points require a separate reduction before
taking derivatives.

\subsection{Continuity and maximum entropy on a face}

For probability laws $p$ and $q$ on the same finite set, write
$\Ent(p)=-\sum_a p_a\log p_a$. We will use the elementary comparison
\begin{equation}\label{ent:comparison}
 \Ent(p)\le-\sum_{a:p_a>0}p_a\log q_a
\end{equation}
whenever $q_a>0$ on the support of $p$; equality holds exactly when $p=q$.
Indeed, $\log t\le t-1$ gives
$\sum_{p_a>0}p_a\log(q_a/p_a)\le\sum_{p_a>0}q_a-1\le0$,
with the asserted equality condition. This also proves
$\Ent(p)\le\log N$ for a law on $N$ atoms, by comparison with the uniform law.

\begin{lemma}[Continuity at the boundary]\label{ent:continuity}
For a nonempty perfect-matching polytope $P(G)$, the maximum mean entropy
$H$ is attained, nonnegative, concave and continuous on the whole polytope.
\end{lemma}
\begin{proof}
The laws of a prescribed mean form a nonempty closed subset of the finite
probability simplex. Compactness and continuity of entropy give attainment.
Mixing maximizing laws proves concavity, and entropy is nonnegative.

For upper semicontinuity at $x$, let $y_j\to x$ in $P(G)$ and choose
maximizing laws at $y_j$. Take a subsequence realizing the upper limit of
their entropies and then a convergent subsequence of laws. Its limit has
mean $x$, so continuity of entropy on the simplex gives
$\limsup_j H(y_j)\le H(x)$.

For lower semicontinuity, take $y\ne x$ close to $x$ in $P(G)$ and put
$t=\sqrt{\|y-x\|_2}$ and $z=x+(y-x)/t$. For sufficiently close $y$,
we have $0<t<1$ and $z\in P(G)$. To check the latter assertion, use
the finite linear description of the polytope. Equalities are preserved;
an inequality tight at $x$ is preserved because its slack at $z$ is its
slack at $y$ divided by $t$; and every inactive inequality stays satisfied
since $z\to x$. There are only finitely many inactive inequalities.
Since $y=(1-t)x+tz$, concavity and nonnegativity give
$H(y)\ge(1-t)H(x)$. Letting $y\to x$ proves the reverse semicontinuity.
\end{proof}

Let $\mathcal A$ be the set of matching incidence vectors, let $F_x$ be
the minimal face containing $x$, and put $\mathcal A_x=\mathcal A\cap F_x$.
The next lemma is the finite maximum-entropy principle on that face.
The exponential representation is the standard finite maximum-entropy
duality; compare Singh and Vishnoi
\citep[full version, Definition~2.1 and Lemma~2.3]{SinghVishnoi2014}.
Jaynes~\citep{Jaynes1957} gives the historical maximum-entropy principle,
and the covariance identities are standard in exponential-family theory
\citep[Proposition~3.1]{WainwrightJordan2008}. We work on the minimal face so that the
relative-interior condition holds, and include the support and boundary
arguments needed here.

\begin{lemma}[Exponential laws on the minimal face]\label{ent:exponential}
Every law with mean $x\in P(G)$ is supported on $\mathcal A_x$.
Its entropy maximizer is positive on every atom of $\mathcal A_x$ and has
the form
\begin{equation}\label{ent:gibbs}
 p_a=\frac{e^{\theta\cdot a}}{\sum_{b\in\mathcal A_x}e^{\theta\cdot b}}
 \quad(a\in\mathcal A_x)
\end{equation}
for a finite vector $\theta\in\R^E$.
For every $\eta\in\R^E$, the law obtained by replacing $\theta$ by $\eta$
maximizes entropy at its own mean, and that mean has minimal face $F_x$.
\end{lemma}
\begin{proof}
The active inequalities defining $F_x$ have nonnegative slack at each
matching vector. Their expected slack is zero at $x$, so every
positive-probability atom belongs to $F_x$. The geometric description also
shows that $F_x=\conv\mathcal A_x$ and that $x$ is relatively interior to it.

Let $b$ be the barycenter of $\mathcal A_x$. For sufficiently small
$\varepsilon>0$, the point $(x-\varepsilon b)/(1-\varepsilon)$ belongs
to $F_x$. Represent it by a law on $\mathcal A_x$ and mix in weight
$\varepsilon$ of the uniform law. This produces a law $q$ of mean $x$
positive on every face atom. If a maximizing law $p$ vanished on a
nonempty set $Z\subseteq\mathcal A_x$, then
\[
 \Ent((1-u)p+uq)-\Ent(p)
 =u\log(1/u)\sum_{a\in Z}q_a+O(u)>0
\]
for sufficiently small $u>0$. The summands outside $Z$ are differentiable
at $p_a>0$, which justifies the $O(u)$ term. Thus $p$ is positive.

Form the matrix with columns $(1,a)$, indexed by $a\in\mathcal A_x$.
All sufficiently small perturbations of $p$ in its kernel preserve
positivity, mass and mean. The entropy gradient $(-1-\log p_a)_a$
therefore lies in the row space. Because the constant row is present,
$\log p_a=c+\theta\cdot a$ for finite real coefficients; normalization
gives Equation~\eqref{ent:gibbs}. No independence of the constraint rows
is required.

For a law of the displayed form with parameter $\eta$, all face atoms
have positive probability. An inequality is tight at its mean $y$
exactly when it is tight at every atom of $\mathcal A_x$, equivalently
when it is tight at $x$. Hence $F_y=F_x$, and every law $q'$ of mean $y$
is supported on $\mathcal A_x$. Equation~\eqref{ent:comparison} gives
\[
 \Ent(q')\le\log\sum_{a\in\mathcal A_x}e^{\eta\cdot a}-\eta\cdot y,
\]
with equality for the exponential law itself. This proves maximality.
\end{proof}

\subsection{The trace obstruction}

For a real number $c$, define $J_c=F-cB$. The following lemma isolates
the analytic consequence of a maximum; it does not yet choose $c$.

\begin{lemma}[Second variation at a fractional maximum]\label{ent:trace}
Let $G$ have $2m\ge2$ vertices, and suppose that $x\in P(G)$ has
$0<x_e<1$ for every edge. If $J_c-H$ has a local maximum at $x$, then,
with $d=\dim F_x$ and $s=|E|$,
\begin{equation}\label{ent:trace-obstruction}
                         0\ge d-s+(c+1)m.
\end{equation}
\end{lemma}
\begin{proof}
Let $p$ and $\theta$ be supplied by Lemma~\ref{ent:exponential}, and
write $T=\operatorname{span}(F_x-x)$. For $v\in\R^E$ and real $t$, set
\[
 \Phi(t)=\log\sum_{a\in\mathcal A_x}e^{(\theta+tv)\cdot a},\qquad
 p_a(t)=e^{(\theta+tv)\cdot a-\Phi(t)},\qquad
 y(t)=\sum_a p_a(t)a.
\]
The lemma gives $F_{y(t)}=F_x$ and
\begin{equation}\label{ent:tilt-entropy}
              H(y(t))=\Phi(t)-(\theta+tv)\cdot y(t).
\end{equation}
All sums are finite, so these curves are smooth. Every coordinate of
$y(t)$ remains strictly between zero and one, since both coordinate
values occur among the positive-probability face atoms.

Define the covariance matrix
\[
 K=\sum_{a\in\mathcal A_x}p_a(a-x)(a-x)^{\mathsf T}.
\]
Differentiation gives $p_a'(0)=p_a v\cdot(a-x)$ and $y'(0)=Kv$.
For every $w$, we have
$w^{\mathsf T}Kw=\sum_a p_a(w\cdot(a-x))^2$.
Positivity of $p$ and the spanning of $T$ by $a-x$ therefore give
\begin{equation}\label{ent:covariance-range}
                 \ker K=T^\perp,\qquad \operatorname{im}K=T,
                 \qquad \rank K=d.
\end{equation}
Also $y''(0)\in T$, since the curve stays in $x+T$.

Let $g=\nabla J_c(x)$ and let $D$ be its diagonal Hessian:
\[
 g_e=-\log x_e-c\log(1-x_e)-(c+1),\qquad
 D_{ee}=-\frac1{x_e}+\frac{c}{1-x_e}.
\]
Since $\Phi'(t)=v\cdot y(t)$, Equation~\eqref{ent:tilt-entropy}
gives $\frac{d}{dt}H(y(t))=-(\theta+tv)\cdot y'(t)$.
For $r(t)=J_c(y(t))-H(y(t))$, the first two derivatives at zero are
\[
 \begin{aligned}
 r'(0)&=(g+\theta)\cdot Kv,\\
 r''(0)&=(Kv)^{\mathsf T}D(Kv)+v^{\mathsf T}Kv
                       +(g+\theta)\cdot y''(0).
 \end{aligned}
\]
Stationarity gives $r'(0)=0$ for every $v$. Since $\operatorname{im}K=T$
and $y''(0)\in T$, the last term in $r''(0)$ vanishes. The second
derivative test is therefore
\begin{equation}\label{ent:second-variation}
 0\ge (Kv)^{\mathsf T}D(Kv)+v^{\mathsf T}Kv
 \qquad(v\in\R^E).
\end{equation}
This cancellation is why differentiating within the minimal face suffices;
no smoothness of $H$ across different faces has been assumed.

Take orthonormal eigenvectors $u_1,\ldots,u_d$ of $K$ with positive
eigenvalues $\lambda_1,\ldots,\lambda_d$. Substituting $v=u_i$ in
Equation~\eqref{ent:second-variation}, dividing by $\lambda_i$, and
summing gives
\[
             0\ge d+\sum_i\lambda_i u_i^{\mathsf T}Du_i
                    =d+\tr(KD).
\]
Only the spectral decomposition of $K$ is used; $K$ and $D$ need not
commute. If $d=0$, this same statement is the empty sum identity.
Each coordinate is Boolean, so $K_{ee}=x_e(1-x_e)$. Using
$\sum_e x_e=m$, we obtain
\[
 \tr(KD)=\sum_e\bigl[-(1-x_e)+cx_e\bigr]=-s+(c+1)m,
\]
which proves Equation~\eqref{ent:trace-obstruction}.
\end{proof}

\begin{proof}[Lower bound in Theorem~\ref{ent:main}]
Write the original vertex count as $2M\ge2$, and fix
$c=2-2/M+\varepsilon$ with $\varepsilon>0$.
By Lemma~\ref{ent:continuity}, $J_c-H$ attains a maximum on $P(G)$.
Suppose that this maximum is positive, at $x$. Delete its zero edges
and the endpoints of its value-one edges. The exact fractional reduction
in Lemma~\ref{geom:lem-reduction} identifies all laws of mean $x$ with laws
of the reduced mean and preserves $F,B,H$. An empty reduction would give
$F=B=H=0$, contrary to positivity. The reduced graph therefore has
$2m$ vertices for some $1\le m\le M$, and all retained coordinates are
strictly fractional.

Every feasible mean in the reduced graph extends by the fixed matching
edges to a feasible original mean, preserving the same three functions.
The reduced mean consequently remains a global maximum of $J_c-H$.
The geometric rank bound gives $d\ge s-3m+2$, whereas
Lemma~\ref{ent:trace} gives
\[
 0\ge d-s+(c+1)m\ge 2+(c-2)m
       =2\left(1-\frac mM\right)+\varepsilon m>0.
\]
This contradiction proves $H\ge F-cB$ for every $\varepsilon>0$.
For each fixed mean, $B$ is finite; letting $\varepsilon\downarrow0$
proves $H\ge F-(2-2/M)B$. The strict coefficient followed by this limit
is essential: the displayed geometric comparison can be an equality
when $m=M$ and $\varepsilon=0$.
\end{proof}

\subsection{The upper bound and componentwise refinement}

The upper bound is the vertex-star instance of Shearer's entropy
inequality~\citep[Section~V, inequality~(22)]{ChungGrahamFranklShearer1986}.
We give the short chain-rule proof to make the counting of parallel
edge coordinates explicit.

\begin{proof}[Upper bound in Theorem~\ref{ent:main}]
Let $X=(X_e)_{e\in E}$ be the incidence vector of any matching law with
mean $x$, and order $E$. For a vertex $v$, the chain rule gives
\[
 \Ent(X_{\delta(v)})
 =\sum_{e\in\delta(v)}
   \Ent\bigl(X_e\mid X_f:f\in\delta(v),\ f<e\bigr)
 \ge\sum_{e\in\delta(v)}\Ent\bigl(X_e\mid X_f:f<e\bigr).
\]
Conditioning on additional finite random variables decreases conditional
entropy, by concavity of entropy applied to the conditional laws.
Every edge, including every labelled parallel edge, belongs to exactly
two vertex stars. Summing and applying the chain rule again therefore gives
$\sum_v\Ent(X_{\delta(v)})\ge2\Ent(X)$.
At a vertex the star vector makes exactly one choice, with probabilities
$(x_e)_{e\in\delta(v)}$, so
\[
 \Ent(X)\le\frac12\sum_v\left[-\sum_{e\in\delta(v)}x_e\log x_e\right]
 =F(x).
\]
Taking the maximum over laws proves the upper bound, including fixed
coordinates and disconnected graphs.
\end{proof}

\begin{corollary}[Fractional components]\label{ent:components}
For $x\in P(G)$, remove its zero edges and the endpoints of its value-one
edges. Let the nonempty connected components of the remaining graph have
$2m_i$ vertices, and let $B_i$ denote the sum defining $B$ on component $i$.
Then
\[
                 F(x)-\sum_i(2-2/m_i)B_i(x)\le H(x)\le F(x).
\]
If there are no remaining components, $F(x)=B(x)=H(x)=0$.
\end{corollary}
\begin{proof}
The exact reduction preserves the functions. Every residual component
has a perfect matching, hence even positive order. A law of a whole
matching induces matching laws on the components, whose entropies sum
to at least the joint entropy. Conversely, the independent product of
component maximizing laws realizes the prescribed mean and the sum of
their entropies. Thus $H$, like $F$ and $B$, is additive at these fixed
component means. Apply Theorem~\ref{ent:main} separately to each component.
\end{proof}
